% ============================================================
\section{Nền tảng và Nghiên cứu liên quan}
\label{sec:related}
% ============================================================

\subsection{Tóm tắt văn bản trích xuất dựa trên đồ thị}
Tóm tắt văn bản tự động rút gọn văn bản gốc nhưng vẫn bảo toàn thông tin cốt
lõi. So với tóm tắt tóm lược (abstractive), tóm tắt trích
xuất (extractive) có ưu thế về tốc độ, tính trung thực và không cần dữ liệu huấn
luyện -- những yếu tố then chốt trong bối cảnh không giám sát và tránh ``ảo
giác'' (hallucination). Trong họ
trích xuất, hướng dựa trên đồ thị mô hình hóa văn bản thành mạng lưới câu và đánh
giá tầm quan trọng dựa trên cấu trúc toàn cục, không cần huấn luyện. TextRank áp
dụng PageRank vào miền văn bản và được
dùng rộng rãi~\cite{mihalcea2004textrank}, nhưng chất lượng phụ thuộc vào hàm đo
độ tương đồng. Phiên bản gốc dùng Word Overlap chỉ phản ánh tương đồng bề mặt từ
vựng; nhiều cải tiến đã được đề xuất (TF-IDF, LexRank, BM25+, và TextRank mở
rộng kết hợp LDA với hệ số giảm chấn điều chỉnh động~\cite{vora2024extractive}) song
vẫn giới hạn ở đặc trưng từ vựng.
Các nghiên cứu khuyến nghị chuyển sang sentence embedding để nắm bắt ngữ nghĩa
sâu hơn~\cite{kadriu2021extractive}.

\subsection{Trích xuất cụm từ khóa}
Các phương pháp trích xuất keyphrase không giám sát chia thành ba
nhóm. Nhóm \emph{thống kê} (TF-IDF, RAKE~\cite{rose2010automatic},
YAKE~\cite{campos2020yake}) nhanh nhẹ nhưng không
nắm bắt ngữ nghĩa. Nhóm \emph{đồ thị} (TextRank~\cite{mihalcea2004textrank},
TopicRank~\cite{bougouin2013topicrank}, Multipartite~\cite{boudin2018unsupervised},
WordAttractionRank~\cite{wang2014corpus}) cải thiện hơn nhưng thiếu cơ chế giảm
trùng lặp ở mức ngữ nghĩa. Nhóm \emph{embedding} đưa cả văn bản và ứng viên vào
cùng không gian vector: KeyBERT~\cite{grootendorst2020keybert} chọn các n-gram có
độ tương đồng cosine cao nhất với văn bản và có thể tích hợp Maximal Marginal
Relevance (MMR)~\cite{carbonell1998use} để tăng đa dạng. Một điểm mấu chốt là
\emph{tính bất đối xứng} của bài toán KeyBERT: một bên là tài liệu dài, một bên
là keyphrase rất ngắn; nếu mô hình embedding không được huấn luyện cho quan hệ
``truy vấn ngắn -- tài liệu dài'', chất lượng sẽ không tối ưu.

\subsection{Khám phá không gian nhãn}
Mô hình hóa chủ đề trải qua ba thế hệ: (i) dựa trên ``túi từ'' như LDA, NMF;
(ii) dựa trên embedding và phân cụm như Top2Vec, BERTopic; (iii) dùng LLM để
sinh và đặt tên chủ đề như TopicGPT. Hạn chế chung là thiên về gán
\emph{một} chủ đề cho mỗi văn bản, khó mở rộng sang kịch bản đa nhãn. Gần đây,
X-MLClass~\cite{li2024open} khai thác LLM một cách có hệ thống để khám phá không
gian nhãn đa nhãn (chia chunk, sinh keyphrase, phân cụm, lặp khám phá nhãn hiếm),
đạt độ phủ vượt trội so với các phương pháp truyền thống trên nhiều benchmark
tiếng Anh. Tuy nhiên, X-MLClass đòi hỏi LLM xử lý toàn bộ văn bản thô, dẫn đến
chi phí token lớn, tốc độ chậm, và trên dữ liệu báo chí giàu thực thể dễ sinh
nhãn nhiễu (entity, meta). StatLM kế thừa ý tưởng cốt lõi của X-MLClass nhưng đặt
một encoder thống kê phía trước để lọc nhiễu và chỉ cho LLM làm việc với keyphrase
đã được cô đặc.

\subsection{Tài nguyên xử lý tiếng Việt}
StatLM sử dụng hai công cụ tách từ tùy theo nhu cầu: \emph{pyvi}~\cite{pyvi} (nhẹ,
chỉ tách từ) và \emph{VnCoreNLP}~\cite{vu2018vncorenlp} (đồng thời tách từ, gán
nhãn từ loại và nhận dạng thực thể). Hai mô hình embedding chuyên biệt được dùng:
\emph{dangvantuan/vietnamese-embedding}~\cite{dangvantuan2021vietnamese} -- một
Sentence-BERT~\cite{reimers2019sentence} trên nền PhoBERT~\cite{nguyen2020phobert}
tối ưu cho tương đồng ngữ nghĩa đối xứng (STS, Pearson 84,21 trên
vi-stsbenchmark); và
\emph{bkai-foundation-models/vietnamese-bi-encoder}~\cite{nguyen2025improving} --
một Bi-Encoder huấn luyện cho truy xuất bất đối
xứng, đứng đầu benchmark VCS tiếng Việt.

